% !TeX root = ../main.tex

\section{Assignment \#4}

\begin{problem}[note =
    Molecular Term Symbols and Electronic Band Types in \ce{NO}]
  \ce{NO} is a heteronuclear diatomic molecule. Because it is heteronuclear, its
  molecular term symbols do not contain $g/u$ labels.

  The ground electronic state of \ce{NO} is $X^2\Pi$.

  Several electronically excited states are listed below.
  \begin{center}
    \begin{tabular}{*3l}
      \toprule
      \textbf{State} & \textbf{Molecular term symbol} &
      \textbf{Approximate character}\\
      \midrule
      Ground  state   & $X \prescript2{}\Pi$       &
      One unpaired electron in a $\pi^*$ orbital\\
      Excited state A & $A \prescript2{}\Sigma^+$ & 
      Electronic excitation to a $\sigma$-type state\\
      Excited state B & $B \prescript2{}\Pi$ & 
      Electronic excitation to another $\pi$-type state\\
      Excited state C & $C \prescript4{}\Pi$ & 
      Quartet excited state\\
      \bottomrule
    \end{tabular}
  \end{center}
  For diatomic electronic transitions, the approximate electric-dipole
  selection rules are
  \[
    \Delta S = 0, \quad \Delta \Lambda = 0, \pm 1.
  \]
  A transition with $\Delta \Lambda = 0$ is usually called a
  \emph{parallel band}, while a transition with $\Delta \Lambda = \pm 1$ is
  usually called a \emph{perpendicular band}.
  \paragraph{Tasks}
  \begin{enumext}
    \item Explain why $X^2\Pi$ has two spin-orbit components:
    ${}^2\Pi_{1/2}$ and ${}^2\Pi_{3/2}$.
    \item Determine whether the following transitions are spin-allowed or
    spin-forbidden:
    \begin{enumext}[columns = 3]
      \item $A \prescript2{}\Sigma^+ \leftarrow X \prescript2{}\Pi$
      \item $B \prescript2{}\Sigma^+ \leftarrow X \prescript2{}\Pi$
      \item $C \prescript4{}\Sigma^+ \leftarrow X \prescript2{}\Pi$
    \end{enumext}
    \item Explain why heteronuclear molecules such as \ce{NO} do not use $g/u$
    labels, while homonuclear molecules such as \ce{O2} do.
  \end{enumext}
\end{problem}
\begin{solution}\leavevmode
  \begin{enumext}
    \item A $^2\Pi$ state has orbital angular momentum projection $\Lambda = 1$
    and total spin $S = \frac12$. The projection of the spin onto the
    internuclear axis is $\Sigma = \pm\frac12$. The total angular momentum
    projection quantum number is $\Omega = |\Lambda + \Sigma|$, yielding two
    possible values:
    \[
    \Omega = \ab|1 + \frac12| = \frac32 \qq{and}
    \Omega = \ab|1 - \frac12| = \frac12.
    \]
    Spin-orbit coupling causes states with different $\Omega$ to have different
    energies, so the $^2\Pi$ state splits into two components,
    denoted $^2\Pi_{3/2}$ and $^2\Pi_{1/2}$.
    \item Spin-allowed transitions obey $\Delta S = 0$.
    \begin{enumext}
      \item $A \prescript2{}\Sigma^+ \leftarrow X \prescript2{}\Pi$: both are
      doublets ($S = \frac12$), so $\Delta S = 0$ is spin-allowed.
      \item $B \prescript2{}\Sigma^+ \leftarrow X {^2\Pi}$: both are doublets
      ($S = \frac12$), so $\Delta S = 0$ is spin-allowed.
      \item $C \prescript4{}\Sigma^+ \leftarrow X {^2\Pi}$: the excited state is a
      quartet ($S = \frac32$) while the ground state is a doublet
      ($S = \frac12$), giving $\Delta S = 1$ spin-forbidden.
    \end{enumext}
    \item The labels $g$ (gerade, symmetric) and $u$ (ungerade, antisymmetric)
    describe the parity of a molecular wavefunction with respect to inversion
    through the center of inversion. This symmetry operation exists only in
    molecules with a center of inversion, i.e., homonuclear diatomics like
    \ce{O2}. Heteronuclear diatomics like \ce{NO} lack an inversion center, so
    their electronic states cannot be classified as $g$ or $u$.
  \end{enumext}
\end{solution}

\begin{problem}[note = Vibronic Transitions and the Franck-Condon Principle]
  Consider a molecule whose ground and excited electronic states can both be
  approximated by harmonic potential-energy curves along one normal
  coordinate $Q$. The excited-state potential is displaced relative to the
  ground-state potential, but the vibrational frequency is assumed to be the
  same in both states.
  \begin{remark}[note = Franck-Condon principle]
    Electronic transitions are much faster than nuclear
    motion, so absorption is approximately vertical.
    \begin{equation}\label{eq:Franck-Condon_principle}
      I_{0\to n} \propto \upe^{-S} \frac{S^n}{n!}
    \end{equation}
  \end{remark}
  Given parameter $S = 1.20$.
  \paragraph{Tasks}
  \begin{enumext}
    \item Calculate the relative Franck-Condon intensities for $0 \to 0$,
    $0 \to 1$, $0 \to 2$, $0 \to 3$, and~$0 \to 4$
    using~\eqref{eq:Franck-Condon_principle}.
    \item Identify which vibronic peak is expected to be the strongest.
    \item Explain how the Huang-Rhys factor $S$ reflects the degree of geometry
    change between the ground and excited electronic states.
  \end{enumext}
\end{problem}
\begin{solution}\leavevmode
  \begin{enumext}
    \item Using the given formula
    \[\begin{alignedat}{2}
      I_{0\to0} & = \upe^{-1.2} \frac{1.2^0}{0!} = 0.301,\quad &
      I_{0\to1} & = \upe^{-1.2} \frac{1.2^1}{1!} = 0.361,\\
      I_{0\to2} & = \upe^{-1.2} \frac{1.2^2}{2!} = 0.217,\quad &
      I_{0\to3} & = \upe^{-1.2} \frac{1.2^3}{3!} = 0.087,\quad
      I_{0\to4}   = \upe^{-1.2} \frac{1.2^4}{4!} = 0.026,
    \end{alignedat}\]
    \item The $0 \to 1$ transition has the highest relative intensity
    ($0.361$). Therefore, the $0 \to 1$ vibronic peak is expected to be the
    strongest in the absorption spectrum.
    \item The Huang-Rhys factor $S$ directly quantifies the displacement
    $\Delta Q$ between the equilibrium geometries of the ground and excited
    electronic states along the normal coordinate $Q$. It is defined as
    \[
      S = \frac{\mu\omega}{2\hbar}\,(\Delta Q)^2,
    \]
    where $\mu$ is the reduced mass and $\omega$ the vibrational frequency.
    A larger $S$ corresponds to a larger shift in the equilibrium position,
    meaning a more substantial geometry change upon electronic excitation.
    This, in turn, leads to a broader Franck-Condon envelope with intensity
    spread over higher vibrational quanta. When $S=0$ the two potential minima
    lie directly above each other and only $0\to 0$ has non-zero intensity.
  \end{enumext}
\end{solution}

\begin{problem}[note =
    {Fluorescence, Phosphorescence, and Delayed Fluorescence}]
  After photoexcitation, a molecule can undergo several excited-state processes,
  including fluorescence, phosphorescence, intersystem crossing, and~delayed
  fluorescence. The following photophysical parameters are measured for three
  emitters.\par\noindent
  \begin{minipage}[t]{.48\linewidth}
    \[
      \begin{cases*}
        S   \to S_0 + h\nu_F & (fluorescence)\\
        T_1 \to S_0 + h\nu_p & (phosphorescence)\\
        S_1 \to T_1          & (ISC)\\
        T_1 \to S_1          & (RISC)
      \end{cases*}
    \]
  \end{minipage}
  \hfill
  \begin{minipage}[t]{.48\linewidth}
    \[
      \begin{aligned}
        \tau_F & = \frac1{k_\text{rF} + k_\text{nrS} + k_\text{ISC}}\\
        \Phi_F & =
          \frac{k_\text{rF}}{k_\text{rF} + k_\text{nrS} + k_\text{ISC}}.
      \end{aligned}
    \]
  \end{minipage}\par
  \begin{table}[htbp]
    \centering
    \caption{Photophysical Data}
    \begin{tabularx}{\linewidth}
      {*2{p{.08\linewidth}}p{.16\linewidth}p{.1\linewidth}Xp{.12\linewidth}}
      \toprule
      \textbf{Emitter} & $\boldsymbol{\Delta E_\text{ST}}$  &
      \textbf{Delayed emission spectrum}& \textbf{Lifetime} &
      \textbf{Temperature dependence}   & \textbf{Oxygen sensitivity}\\
      \midrule
      A & \qty{0.75}\eV & red-shifted from fluorescence & \unit\ms-\unit\s  &
      weak                   & strongly quenched\\
      B & \qty{0.06}\eV & same as A                     & \unit\us-\unit\ms &
      stronger at higher $T$ & quenched\\
      C & \qty{0.80}\eV & same as A                     & \unit\us          &
      depends on excitation intensity & quenched\\
      \bottomrule
    \end{tabularx}
  \end{table}
  For Emitter B: $k_\text{rF} = \qty{5.0e7}{\per\s}$,
  $k_\text{nrS} = \qty{1.0e7}{\per\s}$, $k_\text{ISC} = \qty{4.0e7}{\per\s}$.
  \paragraph{Tasks}
  \begin{enumext}
    \item Calculate the prompt fluorescence lifetime $\tau_F$ of Emitter B.
    \item Calculate the prompt fluorescence quantum yield $\Phi_F$ of Emitter B.
    \item Explain why oxygen usually quenches phosphorescence and delayed
    fluorescence more efficiently than prompt fluorescence.
  \end{enumext}
\end{problem}
\begin{solution}\leavevmode
  \begin{enumext}
    \item [(a) (b)] Substituting the parameters into the given equations
    \begin{align*}
      \tau_F(B)
    = \frac{1}{(5 + 1 + 4)\times \num{e7}} = \qty{e-8}{\s},\enspace
      \phi_F(B)
    = \frac{\num{5e7}}{(5 + 1 + 4)\times \num{e7}} = 50\%
    \end{align*}
    \addtocounter{enumXi}{2}
    \item Oxygen preferentially quenches phosphorescence and delayed fluorescence due to two key factors
    \subparagraph{Kinetics (State Lifetimes)}%
    These emissions rely on long-lived
    triplet ($T_1$) states (\unit\ms-\unit\s), giving oxygen ample time to
    diffuse and collide with the excited molecule. Prompt fluorescence
    originates from the singlet ($S_1$) state, which decays too rapidly
    (\unit\ns) for significant collisional quenching to occur.
    \subparagraph{Thermodynamics (Spin Conservation)}%
    Ground-state oxygen is a
    triplet ($\ce{^3O2}$). Quenching a triplet emitter via Dexter energy transfer ($\ce{T1} + \ce{^3O2} \to \ce{S0} + \ce{^1O2}$) is spin-allowed
    and highly efficient because it conserves total spin angular momentum.
  \end{enumext}
\end{solution}